\begin{theorem}[有限 Hankel 核证书与有限窗正规形的共终端主骨架]\label{thm:conclusion-hankel-null-hankelnf-common-skeleton}
固定整数 \(q\ge 2\)，记 \(P_q(\lambda)\in\ZZ[\lambda]\) 为 \(q\)-moment 的首一最小递推特征多项式，\(d_q:=\deg P_q\)，并沿用 \(\mathrm{NullModes}(q)\) 与 \(\mathrm{HANKELNF}_q\) 的记号。
则下列两类有限数据唯一决定同一个主骨架 \((P_q,r_q)\)：
\begin{enumerate}
  \item 任意一组整数 Hankel-null 基
  \[
  \{\alpha^{(1)},\dots,\alpha^{(\Delta(q))}\}\subset \mathrm{NullModes}(q),
  \]
  及其对应多项式
  \[
  A^{(j)}(\lambda):=\sum_k \alpha^{(j)}_k\lambda^k;
  \]
  \item 任意一组由长度 \(2d_q\) 连续窗口与经审计偏移 \(\ell_q\) 构造出的 Hankel-正规形
  \[
  \mathrm{HANKELNF}_q=
  \bigl(d_q,\ell_q;\ H_0^{(\ell_q)},H_1^{(\ell_q)};\ A_q,\alpha,\beta\bigr),
  \qquad
  A_q:=(H_0^{(\ell_q)})^{-1}H_1^{(\ell_q)}.
  \]
\end{enumerate}
更精确地，
\[
\gcd\bigl(A^{(1)}(\lambda),\dots,A^{(\Delta(q))}(\lambda)\bigr)=\pm P_q(\lambda),
\]
而
\[
r_q=\rho(A_q),
\]
并且 \(A_q\) 所承载的最小递推多项式恰为 \(P_q\)。
\end{theorem}

\begin{proof}
由推论 \ref{cor:pom-resonance-gcd-recovers-minpoly}，任一整数 Hankel-null 基对应多项式的最大公因子恰为 \(\pm P_q(\lambda)\)。取首一规范化后，\(P_q\) 遂由有限 Hankel 核证书唯一恢复，从而其 Perron 根 \(r_q\) 亦唯一确定。

另一方面，定义 \ref{def:pom-hankel-normal-form} 与命题 \ref{prop:pom-hankel-realization-protocol} 给出：由长度 \(2d_q\) 的连续窗口可确定可逆 Hankel 块 \(H_0^{(\ell_q)},H_1^{(\ell_q)}\)，进而确定最小实现矩阵
\[
A_q=(H_0^{(\ell_q)})^{-1}H_1^{(\ell_q)}.
\]
推论 \ref{cor:pom-hankel-free-energy-evaluator} 进一步给出
\[
r_q=\rho(A_q),
\]
且 \(A_q\) 的特征值集合正是递推特征多项式的根。由于 \(\mathrm{HANKELNF}_q\) 是由最小 Hankel realization 唯一构造的正规形，其对应首一最小递推多项式即为 \(P_q\)。故两类有限数据共享同一主骨架 \((P_q,r_q)\)。证毕。
\end{proof}

\begin{corollary}[机制同一性的有限窗口塌缩]\label{cor:conclusion-hankel-window-identity-collapses-to-monic-skeleton}
设两条实现路径在同一 \(q\) 下产生序列 \(S_q,S_q'\)，且具有相同 Hankel 秩 \(d_q\)。若它们在某个长度为 \(2d_q\) 的连续窗口上相等，则
\[
\mathrm{HANKELNF}_q(S_q)=\mathrm{HANKELNF}_q(S_q'),
\qquad
P_q(S_q)=P_q(S_q'),
\qquad
r_q(S_q)=r_q(S_q'),
\]
并且
\[
S_q(m)=S_q'(m)\qquad (\forall m\ge 2).
\]
\end{corollary}

\begin{proof}
命题 \ref{prop:pom-moment-finite-sampling-decidable} 已证明：在 Hankel 秩相同的前提下，长度 \(2d_q\) 的连续窗口相等即强制两侧共享同一 \(\mathrm{HANKELNF}_q\)，从而整条序列对所有 \(m\ge 2\) 完全一致。再由定理 \ref{thm:conclusion-hankel-null-hankelnf-common-skeleton}，同一 Hankel-正规形唯一决定同一首一主骨架 \((P_q,r_q)\)。证毕。
\end{proof}

\begin{theorem}[原始剩余扭结是秩降之外唯一的有限特征失真源]\label{thm:conclusion-hankel-primitive-torsion-only-finite-characteristic-distortion}
在定理 \ref{thm:pom-hankel-primitive-torsion-spectrum} 的记号下，设
\[
K_{\ZZ}:=\ker_{\ZZ}(\mathsf H_{N,L}),
\qquad
M_{\ZZ}(P):=
\Bigl\{
\operatorname{coeff}_{<L}\!\bigl(Q(\lambda)P(\lambda)\bigr):
\deg Q\le L-d-1
\Bigr\},
\]
并定义原始剩余商
\[
T_{\mathrm{prim}}:=K_{\ZZ}/M_{\ZZ}(P).
\]
再记秩降素数集
\[
\mathcal B_{\mathrm{rk}}(\mathsf H_{N,L})
:=
\left\{
p\ \text{prime}:\ \rank_{\FF_p}(\mathsf H_{N,L}\bmod p)<d
\right\}.
\]
则：
\begin{enumerate}
  \item \(T_{\mathrm{prim}}\) 为有限阿贝尔群；
  \item 对每个 \(p\notin \mathcal B_{\mathrm{rk}}(\mathsf H_{N,L})\)，有
  \[
  \dim_{\FF_p}
  \frac{\ker_{\FF_p}(\mathsf H_{N,L}\bmod p)}{\overline M_p}
  =
  \dim_{\FF_p}\!\bigl(T_{\mathrm{prim}}\otimes\FF_p\bigr);
  \]
  \item 特别地，
  \[
  \ker_{\FF_p}(\mathsf H_{N,L}\bmod p)=\overline M_p
  \iff
  p\nmid |T_{\mathrm{prim}}|.
  \]
\end{enumerate}
因此，在排除秩降素数之后，全部额外的有限特征 Hankel-null 证书都由且仅由 \(T_{\mathrm{prim}}\) 的 \(p\)-扭结决定。
\end{theorem}

\begin{proof}
这正是定理 \ref{thm:pom-hankel-primitive-torsion-spectrum} 的内容。其证明把包含
\[
M_{\ZZ}(P)\subset K_{\ZZ}
\]
化为 Smith 标准形，从而把模 \(p\) 后出现的额外核维数精确识别为
\[
\dim_{\FF_p}\!\bigl(T_{\mathrm{prim}}\otimes\FF_p\bigr).
\]
因此有限特征世界中的 Hankel 失真严格分裂为两层：一层来自 \(\mathsf H_{N,L}\) 的秩降，另一层来自原始剩余商的扭结。证毕。
\end{proof}

\begin{corollary}[好素数刚性与失真素数集的有限性]\label{cor:conclusion-hankel-good-prime-rigidity-finite-distortion-support}
定义
\[
\mathfrak S_q^{\mathrm{dist}}
:=
\mathcal B_{\mathrm{rk}}(\mathsf H_{N,L})
\cup
\{\,p:\ p\mid |T_{\mathrm{prim}}|\,\}.
\]
则 \(\mathfrak S_q^{\mathrm{dist}}\) 为有限素数集，并且对每个
\[
p\notin \mathfrak S_q^{\mathrm{dist}}
\]
都有严格刚性
\[
\ker_{\FF_p}(\mathsf H_{N,L}\bmod p)=\overline M_p.
\]
换言之，离开 \(\mathfrak S_q^{\mathrm{dist}}\) 后，模 \(p\) Hankel-null 证书不再携带任何新的 primitive 方向，而完全退化为 \(P_q\) 的有限截断倍数。
\end{corollary}

\begin{proof}
由定理 \ref{thm:conclusion-hankel-primitive-torsion-only-finite-characteristic-distortion}，若
\[
p\notin \mathcal B_{\mathrm{rk}}(\mathsf H_{N,L})
\qquad\text{且}\qquad
p\nmid |T_{\mathrm{prim}}|,
\]
则必有 \(\ker_{\FF_p}(\mathsf H_{N,L}\bmod p)=\overline M_p\)。
有限性方面，\(\mathcal B_{\mathrm{rk}}(\mathsf H_{N,L})\) 由某个非零 \(d\times d\) 子式的素因子控制，因此有限；而 \(T_{\mathrm{prim}}\) 为有限阿贝尔群，其素因子支撑也有限。故并集 \(\mathfrak S_q^{\mathrm{dist}}\) 有限。证毕。
\end{proof}

\begin{theorem}[有限 Hankel 失真的单圆—素数账本终对象]\label{thm:conclusion-hankel-single-circle-prime-ledger-terminal-object}
对每个共振阶 \(q\)，令
\[
\mathfrak S_q^{\mathrm{dist}}
:=
\mathcal B_{\mathrm{rk}}(\mathsf H_{N,L})
\cup
\{\,p:\ p\mid |T_{\mathrm{prim}}|\,\},
\qquad
G_q:=\ZZ[(\mathfrak S_q^{\mathrm{dist}})^{-1}].
\]
则其 Pontryagin 对偶存在短正合列
\[
0\longrightarrow \prod_{p\in \mathfrak S_q^{\mathrm{dist}}}\ZZ_p
\longrightarrow \widehat{G_q}
\longrightarrow \TT
\longrightarrow 0.
\]
并且 \(\mathfrak S_q^{\mathrm{dist}}\) 可由核的 \(p\)-主因子完全恢复：
\[
p\in \mathfrak S_q^{\mathrm{dist}}
\iff
\left(\prod_{r\in \mathfrak S_q^{\mathrm{dist}}}\ZZ_r\right)_{(p)}\neq 0.
\]
因此，每个 \(q\)-moment 的全部有限特征 Hankel 失真都可压缩为“一个圆周 + 一个全不连通素数账本核”的终对象。
\end{theorem}

\begin{proof}
由推论 \ref{cor:conclusion-hankel-good-prime-rigidity-finite-distortion-support}，\(\mathfrak S_q^{\mathrm{dist}}\) 是有限素数集。于是可将定理 \ref{thm:conclusion-finite-prime-localization-pontryagin-rigidity} 直接应用于
\[
S=\mathfrak S_q^{\mathrm{dist}},
\]
从而得到短正合列
\[
0\to \prod_{p\in S}\ZZ_p\to \widehat{\ZZ[S^{-1}]}\to \TT\to 0.
\]
代入 \(S=\mathfrak S_q^{\mathrm{dist}}\) 即得所示结果。该定理同时给出：有限素数集 \(S\) 可由核的 \(p\)-主因子是否非平凡完全恢复，于是 \(\mathfrak S_q^{\mathrm{dist}}\) 亦被唯一编码在该全不连通核中。证毕。
\end{proof}
